\section*{第 \ref{ch:g11:hormones-and-reproduction} 章——激素对生殖的调控}

\begin{solution}{exo:g11:hormones-and-reproduction:1}
\omterm{def:g11:hormones-and-reproduction:hormone}{激素}是由腺体分泌入血液、作用于带有其受体的远处\omterm{def:g10:cells-common-unit:cell}{细胞}的分子。\omterm{def:g11:hormones-and-reproduction:hormone}{反馈}是
链条末端产物对链条起始端腺体的作用：
当它减弱指令（维持某一水平恒定）时为负\omterm{def:g11:hormones-and-reproduction:hormone}{反馈}，
当它增强指令（产生激增）时
为正\omterm{def:g11:hormones-and-reproduction:hormone}{反馈}。
\end{solution}

\begin{solution}{exo:g11:hormones-and-reproduction:2}
下丘脑（GnRH）、垂体（LH 和 FSH）、性腺（
睾丸中的睾酮；卵巢中的雌\omterm{def:g11:hormones-and-reproduction:hormone}{激素}和孕\omterm{def:g11:hormones-and-reproduction:hormone}{激素}）。
\end{solution}

\begin{solution}{exo:g11:hormones-and-reproduction:3}
生精小管产生精子；生精小管之间的\omterm{def:g10:cells-common-unit:cell}{细胞}
分泌睾酮。
\end{solution}

\begin{solution}{exo:g11:hormones-and-reproduction:4}
卵泡期（第 1--13 天）：卵泡生长，分泌雌\omterm{def:g11:hormones-and-reproduction:hormone}{激素}。
排卵（第 14 天）：LH 激增。黄体期（第 15--28 天）：黄体
分泌孕\omterm{def:g11:hormones-and-reproduction:hormone}{激素}和雌\omterm{def:g11:hormones-and-reproduction:hormone}{激素}。
\end{solution}

\begin{solution}{exo:g11:hormones-and-reproduction:5}
周期末黄体退化：
孕\omterm{def:g11:hormones-and-reproduction:hormone}{激素}和雌\omterm{def:g11:hormones-and-reproduction:hormone}{激素}下降，使增厚的子宫
内膜崩解并出血。
\end{solution}

\begin{solution}{exo:g11:hormones-and-reproduction:6}
雌\omterm{def:g11:hormones-and-reproduction:hormone}{激素}峰值在第 12 天，LH 峰值在第 13 天，排卵在第 14 天，孕\omterm{def:g11:hormones-and-reproduction:hormone}{激素}
峰值在第 21 天。
\end{solution}

\begin{solution}{exo:g11:hormones-and-reproduction:7}
去势后睾酮被移除，因此其对下丘脑和垂体的抑制作用消失：
LH 升高。注射睾酮则恢复抑制，
LH 随之下降——两个方向上均为负\omterm{def:g11:hormones-and-reproduction:hormone}{反馈}。
\end{solution}

\begin{solution}{exo:g11:hormones-and-reproduction:8}
第 12 天雌\omterm{def:g11:hormones-and-reproduction:hormone}{激素}水平高且持续维持，使\omterm{def:g11:hormones-and-reproduction:hormone}{反馈}转为正\omterm{def:g11:hormones-and-reproduction:hormone}{反馈}，
触发激增。第 20 天雌\omterm{def:g11:hormones-and-reproduction:hormone}{激素}水平适中
而孕\omterm{def:g11:hormones-and-reproduction:hormone}{激素}水平较高，\omterm{def:g11:hormones-and-reproduction:hormone}{反馈}为负\omterm{def:g11:hormones-and-reproduction:hormone}{反馈}：
LH 受到抑制。
\end{solution}

\begin{solution}{exo:g11:hormones-and-reproduction:9}
第 21 天：无论周期多长，黄体期都大约持续 14 天，
因此排卵发生在下次月经前 14 天；
多出的那一周位于卵泡期。
\end{solution}

\begin{solution}{exo:g11:hormones-and-reproduction:10}
避孕药中稳定、适中的雌\omterm{def:g11:hormones-and-reproduction:hormone}{激素}和类孕\omterm{def:g11:hormones-and-reproduction:hormone}{激素}成分
使\omterm{def:g11:hormones-and-reproduction:hormone}{反馈}全月保持为负\omterm{def:g11:hormones-and-reproduction:hormone}{反馈}：FSH 持续过低，
卵泡无法成熟，因此不出现雌\omterm{def:g11:hormones-and-reproduction:hormone}{激素}峰值，没有正\omterm{def:g11:hormones-and-reproduction:hormone}{反馈}，
没有 LH 激增，也没有排卵。
\end{solution}

\begin{solution}{exo:g11:hormones-and-reproduction:11}
外源睾酮抑制他的下丘脑和垂体；
LH 和 FSH 下降；睾丸不再受刺激，停止分泌
睾酮，也停止产生精子，并萎缩。
\end{solution}

\begin{solution}{exo:g11:hormones-and-reproduction:12}
失去雌\omterm{def:g11:hormones-and-reproduction:hormone}{激素}的抑制后，GnRH 和 FSH 升高；卵泡成熟；
其分泌的雌\omterm{def:g11:hormones-and-reproduction:hormone}{激素}峰值（作用于仍能响应的垂体）
触发 LH 激增和排卵。该方法用于排卵失败，
原因是卵泡接受的 FSH 过少。
\end{solution}

\begin{solution}{exo:g11:hormones-and-reproduction:13}
卵泡生长并分泌\omterm{def:g11:hormones-and-reproduction:hormone}{激素}（卵巢与 FSH 的作用正常），但
雌\omterm{def:g11:hormones-and-reproduction:hormone}{激素}峰值未能引发激增：正\omterm{def:g11:hormones-and-reproduction:hormone}{反馈}
在下丘脑或垂体处失效。注射类 LH \omterm{def:g11:hormones-and-reproduction:hormone}{激素}
于峰值时可诱导排卵。
\end{solution}

\begin{solution}{exo:g11:hormones-and-reproduction:14}
黄体若保持存活，会持续分泌孕\omterm{def:g11:hormones-and-reproduction:hormone}{激素}：子宫内膜
得以维持（不发生月经），负\omterm{def:g11:hormones-and-reproduction:hormone}{反馈}使 LH
和 FSH 维持在低水平（不再形成新卵泡，也不排卵）。
该\omterm{def:g11:hormones-and-reproduction:hormone}{激素}在着床后数天内即出现在血液和尿液中，
这正是检测所依据的指标。
\end{solution}

\begin{solution}{exo:g11:hormones-and-reproduction:15}
$12 \times 0.6 \times 0.4 \approx 2.9$：约 3 个胚胎。
三次移植均无活产的概率为 $0.65^3 \approx 0.27$；
至少一次活产的概率约为 73\%。
\end{solution}

\begin{solution}{pb:g11:hormones-and-reproduction:1}
\textbf{1.} 第 14 天：第 13 天的 LH 激增比排卵早一天，
而孕\omterm{def:g11:hormones-and-reproduction:hormone}{激素}只在排卵后出现，
从第 14 天开始升高。

\textbf{2.} 卵泡期为第 1--13 天（共 13 天）；黄体期
为第 14--27 天或第 28 天（约 14 天）。

\textbf{3.} 第 10 天生长中的卵泡；
第 22 天的黄体，雌\omterm{def:g11:hormones-and-reproduction:hormone}{激素}和孕\omterm{def:g11:hormones-and-reproduction:hormone}{激素}均由其分泌。

\textbf{4.} 第 6 天：子宫内膜薄，在雌\omterm{def:g11:hormones-and-reproduction:hormone}{激素}作用下开始增厚。
第 22 天：子宫内膜厚，呈分泌状态，在孕\omterm{def:g11:hormones-and-reproduction:hormone}{激素}作用下为胚胎着床做好准备。
第 28 天：子宫内膜崩解、出血，因两种\omterm{def:g11:hormones-and-reproduction:hormone}{激素}均已下降。

\textbf{5.} 周期之初，上一周期的\omterm{def:g11:hormones-and-reproduction:hormone}{激素}已经下降，
其抑制作用随之消失，因此 FSH 升高；黄体期中，
孕\omterm{def:g11:hormones-and-reproduction:hormone}{激素}和雌\omterm{def:g11:hormones-and-reproduction:hormone}{激素}对其有强烈抑制。

\textbf{6.} 负\omterm{def:g11:hormones-and-reproduction:hormone}{反馈}：卵泡中不断升高的雌\omterm{def:g11:hormones-and-reproduction:hormone}{激素}抑制垂体，
使 FSH 下降。

\textbf{7.} 正\omterm{def:g11:hormones-and-reproduction:hormone}{反馈}：雌\omterm{def:g11:hormones-and-reproduction:hormone}{激素}维持在最高水平约两天，
转而刺激下丘脑和垂体，而非抑制它们，
于是 LH 激增。

\textbf{8.} 负\omterm{def:g11:hormones-and-reproduction:hormone}{反馈}：孕\omterm{def:g11:hormones-and-reproduction:hormone}{激素}与中等水平的雌\omterm{def:g11:hormones-and-reproduction:hormone}{激素}
共同抑制垂体；若雌\omterm{def:g11:hormones-and-reproduction:hormone}{激素}没有持续的峰值，
则无正\omterm{def:g11:hormones-and-reproduction:hormone}{反馈}，也不会出现第二次激增。

\textbf{9.} 孕\omterm{def:g11:hormones-and-reproduction:hormone}{激素}和雌\omterm{def:g11:hormones-and-reproduction:hormone}{激素}骤降；子宫内膜
崩解，并于下一周期第 1 天出血；
FSH 失去抑制后升高。

\textbf{10.} 黄体\omterm{def:g11:hormones-and-reproduction:hormone}{激素}下降，解除了负\omterm{def:g11:hormones-and-reproduction:hormone}{反馈}，
FSH 升高，新的卵泡开始生长：
一个周期的结束，正是下一个周期的信号。

\textbf{11.} 整月中 FSH 和 LH 均处于低水平（由避孕药
中的\omterm{def:g11:hormones-and-reproduction:hormone}{激素}维持负\omterm{def:g11:hormones-and-reproduction:hormone}{反馈}）；没有卵泡成熟。

\textbf{12.} 没有峰值（无卵泡成熟），没有激增，也没有排卵。

\textbf{13.} 避孕药中的\omterm{def:g11:hormones-and-reproduction:hormone}{激素}下降，子宫内膜变薄，
崩解并出血；这是撤退性出血，而不是黄体期的结束
——因为并无黄体。

\textbf{14.} 避孕药中的\omterm{def:g11:hormones-and-reproduction:hormone}{激素}撤去后，负\omterm{def:g11:hormones-and-reproduction:hormone}{反馈}解除，FSH 升高，
卵泡可能开始生长并分泌\omterm{def:g11:hormones-and-reproduction:hormone}{激素}；若其达到雌\omterm{def:g11:hormones-and-reproduction:hormone}{激素}峰值，
则可能出现 LH 激增和排卵；子宫内膜变薄、
宫颈黏液变稠的状态也不再维持。

\textbf{15.} 孕\omterm{def:g11:hormones-and-reproduction:hormone}{激素}单独作用即可充分抑制下丘脑和垂体，
阻止 LH 激增，并使宫颈黏液增稠。

\textbf{16.} 问题在下丘脑或垂体：
卵巢未受刺激。若垂体功能正常，脉冲式注射 GnRH，
或注射 FSH 和 LH，均可恢复卵泡生长和排卵，
从而绕过缺失的环节。

\textbf{17.} 体外受精：从卵巢中采集卵子，
在培养皿中受精，再将胚胎置于子宫内——配子的结合与
输卵管原本承担的运输，此即体外受精的实质。

\textbf{18.} 利用自然的 LH 激增：排卵发生在激增后
约 36 小时，因此应在卵子排出并丢失之前
将其采集。

\textbf{19.} 自然情况下，领先卵泡分泌的雌\omterm{def:g11:hormones-and-reproduction:hormone}{激素}上升，
使 FSH 降低，而较小的卵泡因缺乏 FSH 而退化；
仅最大的卵泡得以存活。注射的 FSH 会打破这种选择，
使所有卵泡都继续生长。

\textbf{20.} 排卵发生在第 14 天；卵泡生长时，负\omterm{def:g11:hormones-and-reproduction:hormone}{反馈}使
FSH 降低，随后正\omterm{def:g11:hormones-and-reproduction:hormone}{反馈}将雌\omterm{def:g11:hormones-and-reproduction:hormone}{激素}峰值转化为
LH 激增；避孕药使\omterm{def:g11:hormones-and-reproduction:hormone}{反馈}整月保持为负，
因此没有卵泡成熟，也没有激增发生。
\end{solution}
